\chapter{Detalii de implementare}

\section{Autentificare și securitate}
Autentificarea utilizatorilor este implementată prin token-uri JWT, fiecare utilizator primind un token de acces la autentificarea cu succes, pe care îl include în headerul \texttt{Authorization} al fiecărei cereri ulterioare adresate serverului. Token-ul conține în payload identitatea utilizatorului și rolul acestuia, serverul putând extrage aceste informații direct din token fără a interoga baza de date la fiecare cerere, ceea ce reduce semnificativ numărul de operații de citire necesare pentru procesarea cererilor autentificate. Token-urile au o durată de valabilitate limitată care, după expirare, obligă utilizatorul să se autentifice din nou pentru a obține un token nou.

Spring Security este configurat să intercepteze toate cererile primite înainte de a ajunge la controller, un filtru dedicat extrăgând token-ul din header, validând semnătura și termenul de valabilitate și populând contextul de securitate cu identitatea utilizatorului. Dacă token-ul lipsește sau nu este valid, cererea este respinsă imediat cu codul 401, fără a mai ajunge la controller.

Parolele utilizatorilor nu sunt stocate în format text în baza de date. Înainte de salvare, fiecare parolă este trecută printr-un algoritm de hashing BCrypt, autentificarea verificând ulterior dacă hash-ul parolei introduse corespunde celui stocat, fără a exista vreodată posibilitatea de a recupera parola originală din baza de date.

Procesul de înregistrare a presupus câteva decizii mai puțin evidente legate de activarea contului. Un cont nou creat rămâne inactiv până la verificarea adresei de email, serverul generând un token unic de verificare trimis automat la adresa specificată la înregistrare, contul activându-se la accesarea linkului care conține token-ul. Pentru utilizatorii minori există un pas suplimentar față de verificarea obișnuită a emailului, platforma solicitând și consimțământul unui părinte sau tutore legal prin trimiterea unui link separat la o adresă specificată la înregistrare, contul rămânând inactiv până la confirmarea din ambele părți. Asta a necesitat gestionarea unor stări intermediare ale contului și verificarea la autentificare că toate etapele de activare au fost finalizate înainte de a permite accesul în platformă (\Cref{fig:login}).

Autorizarea la nivel de endpoint este implementată prin adnotarea \texttt{@PreAuthorize} pe metodele controller-ului, fiecare metodă specificând ce rol trebuie să aibă utilizatorul autentificat pentru a putea apela endpoint-ul respectiv. Un student nu poate accesa funcționalitățile unui profesor chiar dacă cunoaște URL-ul exact al endpoint-ului, serverul respingând cererea cu codul 403 înainte de a executa orice logică (\Cref{fig:security-filter-chain}).

\section{Gestionarea cursurilor și fluxul de aprobare}
Cursurile trec printr-un ciclu de viață definit de la creare până la publicare, statusul unui curs controlând vizibilitatea sa și operațiile permise asupra lui la fiecare moment. Un curs nou are statusul \texttt{DRAFT} și nu este vizibil studenților, profesorul putând edita titlul, descrierea, nivelul de dificultate și tot conținutul din interiorul cursului cât timp acesta se află în această stare. La trimiterea spre publicare, statusul trece în \texttt{IN\_REVIEW}, cursul devenind vizibil administratorului în coada de revizuire, dar în continuare invizibil studenților (\Cref{fig:course-approval}).

Administratorul poate aproba sau respinge cursul. La aprobare statusul devine \texttt{PUBLISHED} și cursul apare imediat în catalog, profesorul primind o notificare automată cu decizia. La respingere statusul revine la \texttt{DRAFT} și profesorul primește o notificare însoțită de observațiile administratorului, putând face modificările necesare și retrimite spre publicare. Practic, un curs poate trece de mai multe ori prin ciclul \texttt{DRAFT} $\rightarrow$ \texttt{IN\_REVIEW} $\rightarrow$ \texttt{DRAFT} înainte de a fi aprobat.

Lecțiile și exercițiile din interiorul unui curs au propriul status de publicare, independent de cel al cursului, ceea ce înseamnă că un profesor poate publica un curs cu doar o parte din lecții vizibile, adăugându-le pe restul ulterior fără a scoate cursul din catalog. Asta a adăugat complexitate interogărilor serverului, care trebuie să filtreze conținutul după mai multe niveluri de status simultan când un student accesează un curs.

Un aspect important al fluxului de aprobare este că un profesor nu poate modifica un curs aflat în \texttt{IN\_REVIEW}, statusul blocând editarea până când administratorul ia o decizie. Această restricție a fost introdusă pentru a evita situația în care profesorul modifică conținutul în timp ce administratorul îl revizuiește, ceea ce ar putea face ca decizia administratorului să nu mai reflecte versiunea curentă a cursului (\Cref{fig:course-review-transitions}).

\section{Sistemul de rezervare a sesiunilor}
Disponibilitatea profesorilor este configurată prin intervale orare recurente săptămânale, fiecare interval specificând ziua săptămânii, ora de start și durata. Pe baza acestor intervale, serverul calculează dinamic sloturile disponibile pentru o perioadă de timp, excluzându-le automat pe cele deja rezervate sau care se suprapun cu alte sesiuni existente. Sloturile nu sunt precalculate și stocate în baza de date, ci generate la cerere, ceea ce înseamnă că disponibilitatea afișată studentului reflectă întotdeauna situația curentă fără riscul de a afișa sloturi expirate (\Cref{fig:availability}).

Procesul de rezervare implică mai mulți pași de validare executați în cadrul aceleiași tranzacții. Serverul verifică că intervalul orar selectat este înca liber, că studentul are suficiente credite pentru tipul de sesiune ales, individual sau de grup, și că nu există o altă rezervare a aceluiași student care se suprapune cu intervalul selectat. Toate aceste verificări și operațiile care urmează, deducerea creditelor și crearea rezervării, sunt grupate într-o singură tranzacție pentru a preveni situațiile în care doi studenți rezervă simultan același slot sau același student face rezervări concurente care îi consumă mai multe credite decât are disponibile (\Cref{fig:booking-validation}).

La confirmarea rezervării, serverul creează automat o cameră în Daily.co prin API-ul acestora, stocând identificatorul camerei alături de datele rezervării. Camera rămâne activă până la finalizarea sau anularea sesiunii, moment în care serverul o șterge automat din Daily.co. Ambii participanți, studentul și profesorul, primesc notificări la confirmarea rezervării și un reminder înainte de sesiune.

La anularea unei rezervări, creditele sunt returnate studentului, cu excepția cazului în care anularea se face cu mai puțin de un anumit interval de timp înaintea sesiunii, serverul aplicând politica de anulare la momentul procesării cererii.

\section{Sistemul de credite și facturare}
Integrarea cu Stripe acoperă întregul flux de facturare, serverul comunicând cu Stripe prin API pentru crearea de abonamente și primind notificări prin webhook pentru evenimentele care necesită acțiuni pe platformă, cum ar fi reînnoirea unui abonament, o plată eșuată sau anularea de către student (\Cref{fig:billing}).

Procesarea webhook-urilor a avut nevoie de mai multă atenție din cauza modului în care Stripe funcționează. Stripe poate retrimite același eveniment de mai multe ori dacă nu primește un răspuns de confirmare în timp util din partea serverului, de exemplu dacă serverul a fost temporar indisponibil sau a răspuns prea lent. Fără un mecanism de protecție, un eveniment de plată retrimis ar putea duce la alocarea dublă a creditelor. Soluția a fost stocarea identificatorului unic al fiecărui eveniment procesat în tabela \texttt{payment\_events}, serverul verificând la fiecare webhook primit dacă ID-ul a mai fost procesat și ignorând evenimentul dacă acesta există deja (\Cref{fig:webhook-idempotency}).

Balanța de credite a fiecărui student este gestionată pintr-un mecanism care previne situația în care două rezervări simultane ale aceluiași student deduc credite din același sold fără să știe una de cealaltă. La fiecare deducere, serverul verifică că soldul nu a fost modificat de o altă operație între momentul citirii și momentul scrierii, reîncercând tranzacția dacă detectează un conflict. Alături de balanța curentă, există și un jurnal de tranzacții care înregistrează fiecare modificare a creditelor cu motivul aferent și momentul de timp, oferind un istoric complet atât pentru depanare, cât și pentru suport.

Studentul poate gestiona singur abonamentul prin Customer Portal-ul Stripe, accesat direct din platformă printr-un link generat de server, putând schimba metoda de plată sau anula abonamentul fără a fi necesară intervenția unui administrator.

\section{Urmărirea progresului}
Progresul unui student este urmărit pe două planuri diferite: lecții și quiz-uri. Fiecare are logica sa proprie de marcare ca finalizat.

O lecție nu poate fi finalizată dacă are exerciții asociate pe care studentul nu le-a rezolvat corect, serverul verificând la cererea de marcare a progresului că toate exercițiile obligatorii au răspunsuri corecte înregistrate în baza de date. Această decizie a fost luată pentru a evita situația în care un student navighează prin lecții fără a interacționa efectiv cu conținutul, progresul urmărit reflectând astfel o parcurgere reală, nu doar deschiderea paginilor. Lecțiile fără exerciții pot fi marcate ca fiind finalizate imediat, studentul confirmând că a citit conținutul (\Cref{fig:progress}).

Quiz-urile sunt urmărite separat față de lecții, fiecare susținere creând o înregistrare nouă cu răspunsurile date de student și cu rezultatul calculat de server la momentul trimiterii. Limita de timp a unui quiz este verificată exclusiv pe server, nu pe client, prevenind depășirea ei prin manipularea interfeței sau prin trimiterea manuală a cererilor HTTP. Dacă la momentul primirii răspunsurilor serverul detectează că limita de timp a fost depășită, quiz-ul este marcat automat ca fiind picat indiferent de răspunsurile trimise, un exemplu putând fi atunci când un quiz are o limită de 10 minute și răspunsurile sunt trimise la 12 minute de la începere. Studentul poate vedea istoricul tuturor susținerilor unui quiz, inclusiv răspunsurile date și scorul obținut la fiecare.

Procentul de completare al unui curs este calculat pe baza numărului de lecții finalizate față de numărul total de lecții publicate, actualizându-se automat la fiecare lecție terminată. Un curs este considerat complet atunci când studentul a terminat toate lecțiile publicate din acel curs, serverul marcând automat cursul ca fiind terminat fără a fi necesară o acțiune explicită din partea studentului (\Cref{fig:lesson-completion-gate}).

\section{Sistemul de notificări}
Notificările reprezintă un mecanism prezent în aproape toate modulele platformei, fiecare acțiune importantă a unui utilizator generând automat o notificare pentru cel implicat, fie că este vorba de rezervarea unei sesiuni, aprobarea unui curs sau respingerea cererii unui profesor (\Cref{fig:notifications}).

Fiecare notificare este stocată în baza de date cu informații despre destinatar, tipul evenimentului, mesajul și dacă a fost citită, utilizatorul putând accesa istoricul oricând, nu doar în momentul în care notificările au fost generate, ceea ce e util mai ales dacă nu a intrat pe platformă timp de câteva zile și vrea să vadă dintr-o privire ce s-a întâmplat, fără să fie nevoie să verifice manual fiecare secțiune. Practic, orice flux important din platformă generează o notificare, profesorul primind una când administratorul decide ceva legat de cererea sau de cursul lui, studenții primind confirmare la rezervare, iar toți participanții primind un reminder înainte de sesiunile programate.

Notificările sunt create în stratul „service” în aceeași tranzacție cu operația care le declanșează, garantând că, dacă operația eșuează, notificarea nu se creează, evitând situațiile în care un profesor primește confirmare de aprobare, dar cursul lui rămâne în așteptare în baza de date (\Cref{fig:notification-service}).

Nu e cel mai complex modul din platformă, dar fără acesta experiența utilizatorilor ar fi considerabil mai dificilă.
